\section{Literature Review}

A Bitcoin transaction is a variable-sized data structure with a theoretical maximum size of 1 MB.
It consists of multiple fields, including version information, inputs, outputs, and locking scripts.
The \texttt{scriptPubKey} field in transaction outputs defines the spending conditions and determines the transaction type.

The transaction format also enables data embedding and cryptographic commitments through specific output constructions.
This capability has been widely used as a basis for information anchoring, enabling verifiable integrity guarantees without modifying the underlying protocol.

\subsection{Data Anchoring Techniques}

Although Bitcoin was originally designed for peer-to-peer monetary transfers, its globally replicated and immutable ledger motivated its use for applications beyond payments.
Lemieux~\cite{Lemieux2016Trusting} evaluates blockchain technology as a mechanism for preserving the integrity of digital records under various ISO standards.
The study concludes that while blockchain-based systems can provide short- and medium-term integrity guarantees, they do not ensure initial record reliability and present limitations as long-term preservation mechanisms.
In particular, the analysis highlights the long-term reliance on cryptographic primitives such as hash functions and digital signatures, including SHA-256 and RIPEMD-160, as a fundamental concern.

Except for coinbase transactions, Bitcoin does not natively provide a designated field for arbitrary metadata.
Gao and Nobuhara~\cite{Gao2017Decentralized} propose embedding information into standard P2PKH transactions by encoding metadata within so-called fake addresses.
By creating a transaction with $N$ outputs, this approach allows the insertion of up to $20 \times N$ bytes.
However, the method permanently destroys the associated Bitcoin by rendering the outputs unspendable, resulting in high economic cost and limited scalability.

Sward et al.~\cite{Sward2018DataInsertion} conduct a comprehensive analysis of data insertion techniques in Bitcoin transactions.
Their study categorizes approaches into ScriptHash-based constructions and fake address techniques.
While ScriptHash-based approaches are not constrained by a fixed data size limit, they require non-standard scripts and increased construction complexity.
The authors conclude that no universally optimal data insertion method exists, as all approaches involve trade-offs between data capacity, transaction cost, network health, and standardness.

Bartoletti et al.~\cite{Bartoletti2019BitcoinMetadata} further analyze metadata embedding techniques and their impact on Bitcoin’s primary function.
They identify Unspent Transaction Output (UTXO) set growth as a critical concern, as excessive UTXO expansion negatively affects node performance.
Among the evaluated methods, approaches based on \texttt{OP\_RETURN} minimize UTXO bloat by explicitly marking outputs as provably unspendable.

The \texttt{OP\_RETURN} opcode enables arbitrary data insertion by terminating script execution.
Strehle and Steinmetz~\cite{Strehle2020DominatingOPReturns} analyze over 32 million transactions and demonstrate that \texttt{OP\_RETURN} has become the dominant mechanism for embedding non-monetary data in Bitcoin.
The study shows that most such transactions originate from a small number of services, notably Omni and Veriblock.
Widely deployed protocols such as Blockcerts~\cite{Blockcerts} and OpenTimestamps~\cite{OpenTimestamps} rely on \texttt{OP\_RETURN}-based anchoring.
Despite its efficiency and standardization, \texttt{OP\_RETURN} explicitly reveals the presence and location of embedded metadata, limiting privacy for sensitive verification use cases.

Legacy data anchoring techniques therefore exhibit two primary limitations: lack of privacy and transaction malleability.
While SegWit effectively resolves malleability, legacy approaches still expose anchored data at the protocol level.
Kedziora et al.~\cite{Kedziora2022AnalysisSegWit} analyze the impact of SegWit adoption, showing that segregating witness data improves both transaction efficiency and security.
They report transaction cost reductions of up to 75\% and expanded data capacity through the Witness structure without increasing UTXO bloat.

Table~\ref{tab:comparativeBitcoin} summarizes the main characteristics of Bitcoin-based data anchoring techniques discussed in the literature.
The values are approximate and derived from the referenced works, providing a qualitative comparison of trade-offs.

\begin{table}[h]
    \centering
    \resizebox{\textwidth}{!}{%
        \begin{tabular}{|l|c|c|c|c|c|}
            \hline
            \textbf{Characteristic} & \textbf{Fake Address} & \textbf{OP\_RETURN} & \textbf{OpenTimestamps} & \textbf{Script Methods} & \textbf{Witness Methods} \\ \hline
            Size (Bytes) & $20 \times N$ & $\sim$80 & $\sim$40 & $\sim$1000 & $\sim$4000 \\ \hline
            Privacy & No & Hash-visible & Hash & Hidden until spend & Hidden until spend \\ \hline
            Scalability & Limited & Merkle Tree & Merkle Tree & Complex & Complex \\ \hline
            UTXO Impact & High bloat & None & None & Low & Low \\ \hline
            Cost Model & Fee + Burn & Standard & Aggregated & Low & Discounted \\ \hline
        \end{tabular}
    }
    \caption{Comparison of Bitcoin-based data anchoring techniques}
    \label{tab:comparativeBitcoin}
\end{table}

\subsection{Cryptographic Foundations}

Advances in cryptographic primitives further contribute to the evolution of blockchain-based verification systems.
Gayoso et al.~\cite{Gayoso2020AnalysisCryptographic} and Fang et al.~\cite{Fang2020DigitalSignature} survey cryptographic foundations used in blockchain systems, emphasizing the role of hash functions, Merkle trees, and digital signatures.
Yang et al.~\cite{Yang2022EfficientVerifiablyEncrypted} analyze advanced signature constructions, highlighting efficiency and verifiability properties relevant to non-repudiation systems.

More recent work by Zhang et al.~\cite{Zhang2025HighEfficiencySchnorrMultiSignature} demonstrates the efficiency and aggregation advantages of Schnorr-based signatures.
These properties motivate their adoption in Bitcoin through BIP340~\cite{bip340} and their integration into Taproot.

Building upon these developments, BIP341 introduces Taproot as a soft fork extension of SegWit.
Taproot defines a new output type, Pay-to-Taproot (P2TR), combining Schnorr signatures with Merkleized Abstract Syntax Trees (MAST).
Since its activation in November 2021, Taproot has been successfully deployed on the Bitcoin network, enabling more private and efficient spending conditions.
Despite its expressive power, the literature has not yet explored Taproot as the foundation for a complete document verification protocol encompassing integrity, authenticity, and revocation.

\subsection{Alternative Systems and Applications}

To overcome Bitcoin’s script limitations, alternative blockchain platforms have adopted different execution models.
Ethereum replaces the UTXO paradigm with an account-based model and supports a Turing-complete virtual machine.

Numerous studies leverage Ethereum smart contracts for certification and attestation.
Works by Wang et al.~\cite{Wang2019BlockchainCT}, Li et al.~\cite{Li2020DecentralizedPKI}, Pennino et al.~\cite{Pennino2021EfficientCertification}, Palmisano et al.~\cite{Palmisano2022NotarizationAntiPlagiarism}, Mishra et al.~\cite{Mishra2025BlockchainDecentralized}, and Harer et al.~\cite{Harer2022DecentralizedAttestation} demonstrate the flexibility of smart contracts for managing identity, permissions, and document validation.
However, these systems incur increased execution cost, gas fees, and implementation complexity.

Beyond academic credentials, blockchain-based certification has been explored in industrial and supply-chain contexts.
Hawashin et al.~\cite{Hawashin2024UAVTraceability} analyze NFT-based certification for UAV manufacturing components, highlighting scalability and traceability requirements beyond education-focused use cases.

The institutional use of blockchain-based document verification has also been studied extensively.
Castro and Au-Yong-Oliveira~\cite{Castro2021BlockchainHigherEducation} analyze blockchain adoption for higher education diplomas, identifying economic incentives and adoption barriers.
Rustemi et al.~\cite{Rustemi2023SLRBlockchainCertificates} provide a systematic review showing that most academic certificate verification platforms are implemented on Ethereum.
They identify revocation as a major challenge, as modifying immutable smart contracts or credential hashes directly impacts identity consistency.

Recent studies further highlight unresolved challenges related to privacy and revocation.
Ramica et al.~\cite{Ramica2024SelectiveDisclosure} survey selective disclosure mechanisms in digital credentials, while Zhang et al.~\cite{Zhang2025PrivacyProtectionIssuanceRevocationVerifiableCredentials} analyze privacy-preserving revocation schemes in self-sovereign identity systems.
These works demonstrate that privacy-aware revocation remains a complex problem, often requiring additional infrastructure or off-chain coordination.

In summary, the literature reveals a fundamental trade-off.
Bitcoin-based approaches relying on legacy mechanisms such as \texttt{OP\_RETURN} provide low-cost anchoring but suffer from limited privacy and scalability.
Smart contract platforms such as Ethereum enable expressive verification logic but incur higher execution costs and increased complexity.
Recent advances introduced by Taproot suggest that this trade-off may be mitigated at the Bitcoin base layer, motivating further investigation into privacy-preserving and scalable document verification protocols.
